\documentclass[12pt]{article}
\usepackage{fontspec}
\usepackage{polyglossia}
\setdefaultlanguage{russian}
\setmainfont{PT Serif}
\newfontfamily\cyrillicfont{PT Serif}[Script=Cyrillic]
\newfontfamily\cyrillicfonttt{PT Serif}[Script=Cyrillic]
\usepackage{amsmath,amssymb}
\usepackage[margin=1in]{geometry}

\title{Распространеніе закона большихъ чиселъ на величины, зависящія другъ отъ друга\thanks{Transcription from the original: A.\,A.~Markov, ``Extension of the Law of Large Numbers to Quantities Dependent on Each Other,'' \textit{Proceedings of the Physico-Mathematical Society at Kazan University}, 2nd Series, Vol.~15, 1906, pp.~135--156.}}
\author{А.\,А.~Марковъ}
\date{1906}

\begin{document}
\maketitle

% Transcription of: A.A. Markov, "Rasprostranenie zakona bol'shikh chisel
% na velichiny, zavisyashchie drug ot druga", Izvestiya Fiziko-matematicheskogo
% obshchestva pri Kazanskom universitete, 2-ya seriya, tom 15, 1906, ss. 135-156.
% Reprinted in JEHPS, Vol. 2, No. 1b, Novembre 2006.

% --- Page 135 ---

Законъ большихъ чиселъ, въ силу котораго, съ вѣроятностью сколь угодно близкою къ достовѣрности, можно утверждать, что среднее ариѳметическое изъ нѣсколькихъ величинъ, при достаточно большомъ числѣ этихъ величинъ, будетъ произвольно мало отличаться отъ средней ариѳметической изъ ихъ математическихъ ожиданій, выведенъ Чебышевымъ\footnote{Сочиненія П.~Л.~Чебышева. Т.~I. О~среднихъ величинахъ.} изъ разсмотрѣнія математическаго ожиданія квадрата разности между суммой этихъ величинъ и суммой ихъ математическихъ ожиданій. А именно, изъ разсужденій Чебышева ясно, что указанный законъ большихъ чиселъ долженъ оправдываться во всѣхъ тѣхъ случаяхъ, когда математическое ожиданіе квадрата разности между суммой величинъ и суммой ихъ математическихъ ожиданій, при безпредѣльномъ возрастаніи числа величинъ, возрастаетъ медленнѣе чѣмъ квадратъ ихъ числа, такъ что отношеніе этого математическаго ожиданія къ квадрату числа величинъ имѣетъ предѣломъ нуль.

Въ своихъ выводахъ Чебышевъ ограничился простѣйшимъ, и потому наиболѣе интереснымъ случаемъ, независимыхъ величинъ; въ этомъ простѣйшемъ случаѣ, какъ показалъ Че-

% --- Page 136 ---

\noindent бышевъ, математическое ожиданіе вышеуказаннаго квадрата, при безпредѣльномъ возрастаніи числа величинъ, можетъ возрастать только съ такой же быстротой какъ число ихъ, если математическія ожиданія квадратовъ самихъ величинъ остаются конечными, а не возрастаютъ безпредѣльно.

Конечно, условія Чебышева далеко не исчерпываютъ всѣхъ случаевъ, къ которымъ можно примѣнить вышеуказанный законъ, если даже мы ограничимся независимыми величинами. Мы не имѣемъ однако въ виду заниматься разысканіемъ условій необходимыхъ и достаточныхъ для примѣнимости этого закона; а укажемъ только, что выводы Чебышева можно распространить и на нѣкоторые случаи, довольно общаго характера, когда величины зависятъ другъ отъ друга.

\medskip

\noindent\textbf{\S\,1.} На первомъ планѣ можно поставить случай, когда связь величинъ такова, что увеличеніе любой изъ нихъ влечетъ за собой уменьшеніе математическихъ ожиданій остальныхъ.

Остановимся на этомъ случаѣ; пусть разсматриваемыя нами величины будутъ
\[
x_1,\; x_2,\; \ldots\;,\; x_n,\; \ldots\;
\]
а математическія ожиданія ихъ соотвѣтственно равны
\[
a_1,\; a_2,\; \ldots\;,\; a_n,\; \ldots\;;
\]
положимъ еще для краткости
\[
x_l - a_l = z_l.
\]
Разсматривая математическое ожиданіе квадрата
\[
(z_1 + z_2 + \ldots + z_n)^2,
\]
разлагаемъ его на слагаемыя
\[
z_1^2 + z_2^2 + \ldots + z_n^2 + 2z_1 z_2 + 2z_1 z_3 + \ldots\ldots\ldots + 2z_{n-1} z_n
\]

% --- Page 137 ---

\noindent и пользуемся извѣстнымъ предложеніемъ, что математическое ожиданіе суммы равно суммѣ математическихъ ожиданій слагаемыхъ.

Для независимыхъ величинъ математическія ожиданія произведеній вида $z_l z_k$ всѣ равны нулю и потому математическое ожиданіе $(z_1 + z_2 + \ldots + z_n)^2$ приводится къ суммѣ математическихъ ожиданій квадратовъ величинъ $z_1,\; z_2,\; \ldots\;,\; z_n$.

Въ нашемъ же случаѣ математическое ожиданіе каждаго произведенія $z_l z_k$ число отрицательное и потому мат.~ожид.\ $(z_1 + z_2 + \ldots + z_n)^2 {<} \sum$~мат.~ож.\ $z^2_k$ \;$(k = 1, 2, \ldots, n)$.

Чтобы въ этомъ удостовѣриться, положимъ, что совокупность чиселъ
\[
z'_l, z''_l, \;\ldots\;,\; z_l^{(m)},
\]
расположенныхъ въ возрастающемъ порядкѣ, представляетъ всѣ возможныя значенія $z_l$, а вѣроятности ихъ соотвѣтственно равны
\[
p'_l, p''_l, \;\ldots\;, p_l^{(m)};
\]
пусть наконецъ при
\[
z_l = z'_l, z''_l, \;\ldots\;, z_l^{(m)}
\]
математическое ожиданіе $z_k$ соотвѣтственно равно
\[
a'_k, a''_k, \;\ldots\;, a_k^{(m)}.
\]
При такихъ обозначеніяхъ математическое ожиданіе произведенія $z_l z_k$ выражается суммою
\[
p'_l z'_l a'_k + p''_l z''_l a''_k + \ldots + p_l^{(m)} z_l^{(m)} a_k^{(m)},
\]
а математическія ожиданія самихъ величинъ $z_l, z_k$, равныя нулю, могутъ быть представлены въ видѣ суммъ
\[
p'_l z'_l + p''_l z''_l + \ldots + p_l^{(m)} z_l^{(m)}
\]
и
\[
p'_l a'_k + p''_l a''_k + \ldots + p_l^{(m)} a_k^{(m)}
\]

% --- Page 138 ---

\noindent И такъ какъ для разсматриваемаго нами случая, согласно предположенію, должно быть
\[
a'_k > a''_k > \ldots > a_k^{(m)},
\]
то въ силу извѣстнаго неравенства Чебышева\footnote{Сочиненія П.~Л.~Чебышева. Т.~II, стр.~716--719. О~приближенныхъ выраженіяхъ однихъ интеграловъ черезъ другіе. Korkine. \textit{Sur un th\'eor\`eme de M.~Tch\'ebycheff} (Comptes Rendus, T.~XCVI)} имѣемъ
\[
\sum p_l^{(i)}\; z_l^{(i)}\; a_k^{(i)} < \sum p_l^{(i)}\; z_l^{(i)} \sum p_l^{(i)}\; a_k^{(i)} = 0.
\]

Слѣдовательно къ разсматриваемому нами случаю можно примѣнить законъ большихъ чиселъ, если только математическое ожиданіе $z_n^2$ остается конечнымъ, а не возрастаетъ безпредѣльно вмѣстѣ съ~$n$.

Къ тому же неравенству

мат.~ожид.\ $(z_1 + z_2 + \ldots + z_n)^2 < \sum$~мат.~ожид.\ $z^2_k$, а слѣдовательно и къ тому же заключенію о примѣнимости закона большихъ чиселъ, не трудно придти и въ случаѣ когда математическое ожиданіе $x_k$, при всякомъ $k$, уменьшается съ увеличеніемъ суммы
\[
x_1 + x_2 + \ldots + x_{k-1};
\]
для этого нужно только воспользоваться тождествомъ
\begin{align*}
(z_1 + z_2 + \ldots + z_n)^2 = z_1^2 &+ z_2^2 + \ldots + z_n^2 + 2z_1 z_2 \\
&+ 2(z_1 + z_2)z_3 + 2(z_1 + z_2 + z_3)z_4 + \ldots
\end{align*}

Указанныя условія выполняются, напримѣръ, въ случаѣ, когда сумма
\[
x_1 + x_2 + \ldots + x_n
\]
равна числу бѣлыхъ шаровъ среди $n$ шаровъ, вынутыхъ послѣдовательно изъ сосуда при слѣдующихъ условіяхъ:

% --- Page 139 ---

1) первоначальное число бѣлыхъ шаровъ въ сосудѣ равно $2a$, а остальныхъ шаровъ~--- $2b$;

2) вынутые шары обратно въ сосудъ не возвращаются;

3) когда въ сосудѣ остается $a + b$ шаровъ, въ него прибавляютъ $a$ бѣлыхъ шаровъ и $b$ другихъ шаровъ.

Въ силу закона большихъ чиселъ въ данномъ случаѣ, совершенно также какъ въ извѣстномъ случаѣ, когда отношеніе числа бѣлыхъ шаровъ въ числу всѣхъ шаровъ въ сосудѣ сохраняетъ неизмѣнную величину $\dfrac{a}{a+b}$, мы можемъ утверждать, съ вѣроятностью сколь угодно близкою къ достовѣрности, что отношеніе числа появившихся бѣлыхъ шаровъ къ числу всѣхъ вынутыхъ шаровъ, при достаточно большомъ числѣ ихъ, будетъ отличаться отъ $\dfrac{a}{a+b}$ менѣе, чѣмъ на любую заданную величину.

\medskip

\noindent\textbf{\S\,2.} Повторяя, что мы даемъ только достаточныя, но не необходимыя, условія, остановимся на одномъ изъ тѣхъ случаевъ, на которые выводы Чебышева можно распространить по той причинѣ, что вліяніе величинъ
\[
x_1, x_2, \;\ldots\;, x_n, \;\ldots
\]
другъ на друга быстро убываетъ по мѣрѣ увеличенія ихъ взаимнаго разстоянія. Въ нашемъ случаѣ
\[
x_1 + x_2 + \ldots + x_n
\]
будетъ представлять число появленій нѣкотораго событія $A$ при $n$ послѣдовательныхъ испытаніяхъ, связанныхъ такимъ образомъ, что вѣроятность событія $A$ при каждомъ испытаніи имѣетъ вполнѣ опредѣленное значеніе $p'$, если событіе $A$ появилось при непосредственно предшествующемъ испытаніи, и другое опредѣленное значеніе $p''$ въ противномъ случаѣ,

% --- Page 140 ---

\noindent каковы бы ни были результаты прочихъ предшествующихъ ему испытаній; если же результаты всѣхъ испытаній остаются неопредѣленными, то вѣроятность событія $A$ для каждаго изъ нихъ равна одному и тому же числу~$p$.

Приступая къ разсмотрѣнію этого случая, положимъ для краткости
\[
1 - p = q,\; 1 - p' = q',\; 1 - p'' = q''
\]
и замѣтимъ, что числа $p, p', p''$ связаны между собой простымъ равенствомъ
\[
p = pp' + qp'',
\]
на основаніи котораго по даннымъ двумъ изъ этихъ чиселъ не трудно найти третье; оно даетъ
\[
p = \frac{p''}{1 - p' + p''},\; p' = 1 + p'' - \frac{p''}{p},\; p'' = p\,\frac{1 - p'}{1 - p}.
\]

Такъ какъ числа $x_l, x_k$ соотвѣтственно означаютъ число (0 или 1) появленій событія $A$ при испытаніяхъ отмѣченныхъ нумерами $l$ и $k$, то нетрудно установить равенства
\[
a_l = a_k = p.
\]

Нетрудно также убѣдиться, что математическое ожиданіе произведенія $z_l z_k$, равнаго
\[
x_l x_k - a_k x_l + a_l a_k,
\]
приводится къ разности
\[
\text{матем.~ожид.}\; x_l x_k - p^2.
\]

Что же касается математическаго ожиданія произведенія $x_l x_k$, то оно равно вѣроятности появленія событія $A$ при обоихъ испытаніяхъ, отмѣченныхъ нумерами $l$ и $k$, которая въ свою

% --- Page 141 ---

\noindent очередь въ силу теоремы объ умноженіи вѣроятностей можетъ быть выражена произведеніемъ числа $p$, представляющаго вѣроятность появленія событія при испытаніи съ нумеромъ $l$, на нѣкоторое число $R^l_k$, представляющее вѣроятность появленія событія $A$ при испытаніи съ нумеромъ $k$, когда извѣстно, что событіе $A$ появилось при испытаніи съ нумеромъ~$l$.

Придя такимъ образомъ къ равенству
\[
\text{матем.~ожид.}\; z_l z_k = p(R^l_k - p),
\]
мы должны заняться разысканіемъ числа $R^l_k$, которое зависитъ, какъ не трудно убѣдиться, только отъ разности $k - l$ и потому можетъ быть обозначено болѣе простымъ символомъ
\[
R_{k-l}.
\]

При $k - l = 1$ въ силу нашихъ условій имѣемъ
\[
R_1 = p';
\]
затѣмъ нетрудно послѣдовательно найти
\begin{align*}
R_2 &= R_1 p' + (1 - R_1) p'' = p'p' + q'p'', \\
R_3 &= R_2 p' + (1 - R_2) p'', \\
&\cdots\cdots\cdots\cdots\cdots \\
R_{m+1} &= R_m p' + (1 - R_m) p'' = p'' + (p' - p'')\, R_m.
\end{align*}
Уравненіе
\[
R_{m+1} = p'' + (p' - p'')\, R_m
\]
принадлежитъ къ числу тѣхъ, для рѣшенія которыхъ нетрудно указать общую формулу; а именно, это уравненіе даетъ намъ
\[
R_m = p + C(p' - p'')^m,
\]
при чемъ постоянное число $C$ опредѣляется изъ условія
\[
R_1 = p'
\]
и оказывается равнымъ $1 - p = q$.

% --- Page 142 ---

Итакъ
\[
\text{матем.\ ожид.}\ z_l z_k = pq(p' - p'')^{k-l}
\]
и слѣдовательно
\begin{align*}
\text{матем.\ ожид.}\ (z_1 + z_2 + \ldots + z_{k-1})z_k &= \\
= pq[p' - p'' + (p' - p'')^2 + \ldots + (p' - p'')^{k-1}]&.
\end{align*}

Отсюда слѣдуетъ, что при $p' < p''$ математическое ожиданіе произведенія
\[
(z_1 + z_2 + \ldots + z_{k-1})z_k
\]
число отрицательное и потому
\[
\text{матем.\ ожид.}\ (z_1 + z_2 + \ldots + z_n)^2 < npq;
\]
если же $p' > p''$, то математическое ожиданіе произведенія
\[
(z_1 + z_2 + \ldots + z_{k-1})z_k
\]
меньше
\[
\frac{pq(p' - p'')}{1 - p' + p''}
\]
и потому
\[
\text{мат.\ ожд.}\ (z_1 + z_2 + \ldots + z_n)^2 < npq\!\left(1 + \frac{2(p' - p'')}{1 - p' + p''}\right)
< \frac{npq\,(1 + p' - p'')}{1 - p' + p''}.
\]

Такимъ образомъ мы пришли къ неравенствамъ, которыя ясно обнаруживаютъ, что къ данному случаю можно примѣнить законъ большихъ чиселъ.

Въ силу этого закона мы можемъ, съ вѣроятностью сколь угодно близкою къ достовѣрности, утверждать, что при достаточно большомъ числѣ нашихъ испытаній отношеніе числа появленій событія $A$ къ числу испытаній будетъ разниться отъ $p$ на величину меньшую любого даннаго числа.

\S\,3. Выраженіе математическаго ожиданія произведенія

% --- Page 143 ---

\noindent $z_l z_k$, найденное въ предыдущемъ параграфѣ, можно вывести изъ общихъ формулъ, которыми мы займемся и которыя могутъ служить основаніемъ для дальнѣйшихъ изслѣдованій.

Обозначимъ, при предположеніяхъ предыдущаго параграфа, символомъ $P_{m,k}$ вѣроятность, что въ первыхъ $k$ испытаній событіе $A$ появится ровно $m$ разъ.

Мы можемъ положить
\[
P_{m,k} = V_{m,k} + U_{m,k},
\]
гдѣ $U_{m,k}$ и $V_{m,k}$, означаютъ туже вѣроятность, какъ и $P_{m,k}$, но при слѣдующихъ условіяхъ:

1) $U_{m,k}$~--- при условіи, что при послѣднемъ испытаніи событіе $A$ не имѣетъ мѣста.

2) $V_{m,k}$~--- при условіи, что при послѣднемъ испытаніи событіе $A$ имѣетъ мѣсто.

Въ связи съ величинами $P_{m,k}$, $V_{m,k}$, $U_{m,k}$ введемъ еще три функціи произвольнаго числа $\xi$:
\begin{align*}
\Phi_k &= U_{0,k} + U_{1,k}\xi + U_{2,k}\xi^2 + \ldots + U_{k-1,k}\xi^{k-1}, \\
\Psi_k &= \phantom{U_{0,k} + {}} V_{1,k}\xi + V_{2,k}\xi^2 + \ldots + V_{k,k}\xi^k, \\
\Omega_k &= P_{0,k} + P_{1,k}\xi + P_{2,k}\xi^2 + \ldots + P_{k,k}\xi^k = \Psi_k + \Phi_k.
\end{align*}

Мы покажемъ теперь, что функція $\Omega_k$ можетъ быть опредѣлена какъ коэффиціентъ при $t^k$, въ разложеніи, по возрастающимъ степенямъ произвольнаго числа $t$, довольно простого выраженія.

Для этой цѣли, пользуясь теоремами сложенія и умноженія вѣроятностей, устанавливаемъ два равенства
\begin{align*}
U_{m,k} &= q'\,V_{m,k-1} + q''\,U_{m,k-1}, \\
V_{m,k} &= p'\,V_{m-1,k-1} + p''\,U_{m-1,k-1}
\end{align*}

Примѣняя затѣмъ эти равенства къ функціямъ $\Phi_k$ и $\Psi_k$, получаемъ два уравненія

% --- Page 144 ---

\begin{align*}
\Phi_k &= q'\,\Psi_{k-1} + q''\,\Phi_{k-1}, \\
\Psi_k &= p'\xi\,\Psi_{k-1} + p''\xi\,\Phi_{k-1},
\end{align*}
изъ которыхъ наконецъ посредствомъ исключенія одной изъ двухъ функцій $\Phi$ или $\Psi$, выводимъ
\begin{align*}
\Phi_{k+1} - (p'\xi + q'')\,\Phi_k + (p' - p'')\xi\,\Phi_{k-1} &= 0, \\
\Psi_{k+1} - (p'\xi + q')\,\Psi_k + (p' - p'')\xi\,\Psi_{k-1} &= 0.
\end{align*}

А такъ какъ
\[
\Omega_k = \Phi_k + \Psi_k,
\]
то должно быть также
\[
\Omega_{k+1} - (p'\xi + q'')\,\Omega_k + (p' - p'')\xi\,\Omega_{k-1} = 0.
\]

Отсюда слѣдуетъ, что $\Omega_k$ можно опредѣлить какъ коэффиціентъ при $t^k$ въ разложеніи, по возрастающимъ степенямъ $t$, дроби вида
\[
\frac{A + Bt}{1 - (p'\xi + q'')t + (p' - p'')\xi t^2},
\]
гдѣ $A$ и $B$ не зависятъ отъ $t$.

Для опредѣленія $A$ и $B$ остается разсмотрѣть $\Omega_k$ при двухъ значеніяхъ $k$; при $k = 1$ и $k = 2$ легко находимъ
\begin{align*}
\Omega_1 &= q + p\xi \\
\Omega_2 &= qq'' + (qp'' + p'q')\xi + pp'\xi^2,
\end{align*}
откуда выводимъ
\[
\Omega_0 = 1.
\]

Разлагая же нашу дробь по возрастающимъ степенямъ $t$ и ограничиваясь двумя первыми членами, находимъ
\begin{align*}
A + (B + p'\xi + q'')t &= \Omega_0 + \Omega_1 t \\
&= 1 + (q + p\xi)t
\end{align*}
что даетъ намъ два равенства
\begin{align*}
A &= 1 \\
B &= (p - p')\xi + q - q'',
\end{align*}

% --- Page 145 ---

\noindent изъ которыхъ послѣднее не трудно привести къ такому виду
\[
B = (p'' - p')(q\xi + p)
\]

Итакъ окончательно находимъ
\[
\frac{1 + (p'' - p')(q\xi + p)\,t}{1 - (p'\xi + q'')t + (p' - p'')\xi\,t^2}
= \Omega_0 + \Omega_1 t + \Omega_2 t^2 + \ldots\,.
\]

Эта формула можетъ служить основаніемъ для дальнѣйшихъ изслѣдованій\footnote{См.\ въ Извѣстіяхъ Акад.\ Наукъ за 1907 годъ мою статью <<Изслѣдованіе замѣчательнаго случая зависимыхъ испытаній>>.} и, въ частности, изъ нея нетрудно вывести результаты предыдущаго параграфа, на чемъ однако мы не остановимся.

\S\,4. Для выясненія дѣла приведемъ еще примѣръ, къ которому нельзя примѣнить законъ большихъ чиселъ.

Пусть подобно прежнему (\S\,1) сумма
\[
x_1 + x_2 + \ldots + x_n
\]
равна числу бѣлыхъ шаровъ среди $n$ шаровъ, вынутыхъ послѣдовательно изъ сосуда, только при новыхъ условіяхъ, которыя мы установимъ такъ:

1) первоначальное число бѣлыхъ шаровъ въ сосудѣ равно $a$, а остальныхъ $b$,

2) каждый вынутый изъ сосуда шаръ поступаетъ обратно въ сосудъ вмѣстѣ съ другимъ шаромъ того же цвѣта.

Въ данномъ случаѣ увеличеніе суммы
\[
x_1 + x_2 + \ldots + x_{k-1}
\]
непремѣнно ведетъ къ увеличенію математическаго ожиданія $x_k$ и потому математическое ожиданіе произведенія
\[
(z_1 + z_2 + \ldots + z_{k-1})z_k
\]
оказывается числомъ положительнымъ.

% --- Page 146 ---

Для вычисленія математическаго ожиданія квадрата
\[
(z_1 + z_2 + \ldots + z_n)^2
\]
замѣтимъ, что оно можетъ быть выражено суммою
\[
\sum \!\left(m - n\,\frac{a}{a+b}\right)^{\!2} P^{\;a,b}_{m,n} \quad (m = 0, 1, 2, \ldots, n),
\]
гдѣ $P^{\;a,b}_{m,n}$ означаетъ вѣроятность, что среди $n$ вынутыхъ шаровъ будетъ ровно $m$ бѣлыхъ, и опредѣляется, какъ нетрудно убѣдиться формулою
\[
P^{\;a,b}_{m,n} = \frac{1\cdot 2\cdot 3\cdot\ldots\cdot n \cdot a(a+1)\ldots(a+m-1)\,b(b+1)\ldots(b+n-m-1)}{1\cdot 2\cdot\ldots\cdot m \cdot 1\cdot 2\cdot\ldots\cdot(n-m)(a+b)(a+b+1)\ldots(a+b+n-1)}.
\]

Изъ приведенной формулы вытекаютъ слѣдующія простыя равенства
\begin{align*}
m\,P^{\;a,b}_{m,n} &= \frac{na}{a+b}\,P^{\;a+1,b}_{m-1,n-1} \\[6pt]
m(m-1)\,P^{\;a,b}_{m,n} &= \frac{n(n-1)\,a(a+1)}{(a+b)(a+b+1)}\,P^{\;a+2,b}_{m-2,n-2}
\end{align*}
на основаніи которыхъ получаемъ
\begin{align*}
\sum m\,P^{\;a,b}_{m,n} &= \frac{na}{a+b} \\[6pt]
\sum m(m-1)\,P^{\;a,b}_{m,n} &= \frac{n(n-1)\,a(a+1)}{(a+b)(a+b+1)},
\end{align*}
гдѣ
\[
m = 0, 1, 2, \ldots, n.
\]

Послѣднія равенства, вмѣстѣ съ очевиднымъ равенствомъ
\[
\sum P^{\;a,b}_{m,n} = 1,
\]
даютъ возможность опредѣлить искомую сумму
\[
\sum \!\left(m - n\,\frac{a}{a+b}\right)^{\!2} P^{\;a,b}_{m,n}\,,
\]
для чего ее слѣдуетъ разложить на три части:

% --- Page 147 ---

\[
\sum m(m-1)\,P^{\;a,b}_{m,n} - \left(\frac{2na}{a+b} - 1\right)\!\sum m\,P^{\;a,b}_{m,n}
+ n^2 \!\left(\frac{a}{a+b}\right)^{\!2}\!\sum P^{\;a,b}_{m,n}
\]

Остается произвести простыя выкладки, по выполненіи которыхъ получаемъ
\[
\sum \!\left(m - n\,\frac{a}{a+b}\right)^{\!2} P^{\;a,b}_{m,n} = \frac{nab(n + a + b)}{(a+b)^2(a+b+1)}\,.
\]

Основываясь на этомъ результатѣ, мы можемъ показать, что къ данному случаю законъ большихъ чиселъ не примѣняется; такъ что при достаточно маломъ значеніи $\varepsilon$ вѣроятность неравенствъ
\[
-\varepsilon \leq \frac{m}{n} - \frac{a}{a+b} \leq +\,\varepsilon
\]
не можетъ быть произвольно близка къ единицѣ, какъ бы ни было велико число $n$.

Для этой цѣли обозначимъ буквою $\mathcal{B}$ вѣроятность неисполненія неравенствъ
\[
-\varepsilon \leq \frac{m}{n} - \frac{a}{a+b} \leq +\,\varepsilon
\]
и примемъ во вниманіе, что квадратъ
\[
\left(\frac{m}{n} - \frac{a}{a+b}\right)^{\!2}
\]
не можетъ превосходить наибольшаго изъ чиселъ
\[
\left(\frac{a}{a+b}\right)^{\!2} \text{ и } \left(\frac{b}{a+b}\right)^{\!2},
\]
которое обозначимъ буквою $\xi$.

При такихъ обозначеніяхъ изъ найденной нами формулы нетрудно вывести неравенство

% --- Page 148 ---

\[
\xi\mathcal{B} > \frac{ab}{(a+b)^2(a+b+1)} - \varepsilon^2,
\]
которое показываетъ, что $\mathcal{B}$ не можетъ быть произвольно малымъ, если только
\[
\varepsilon < \sqrt{\frac{ab}{(a+b)^2(a+b+1)}}\,.
\]

Интересно замѣтить, что въ нашемъ примѣрѣ наивѣроятнѣйшая величина числа $m$, которую мы обозначимъ буквою $\mu$, опредѣляется простыми неравенствами
\begin{align*}
(a - 1)n - b + 1 &\leq (a + b - 2)\mu, \\
(a - 1)n + a - 1 &\geq (a + b - 2)\mu,
\end{align*}
и отношеніе $\dfrac{\mu}{n}$, при безпредѣльномъ возрастаніи $n$, имѣетъ предѣломъ

не $\dfrac{a}{a+b}$\quad а\quad $\dfrac{a - 1}{a + b - 2}$\,;

но конечно нельзя расчитывать, чтобы при большихъ значеніяхъ $n$ отношеніе $\dfrac{m}{n}$ было произвольно близко къ $\dfrac{a - 1}{a + b - 2}$\,;

такъ какъ математическое ожиданіе квадрата $\left(\dfrac{m}{n} - g\right)^{\!2}$ достигаетъ своей наименьшей величины при $g = \dfrac{a}{a+b}$\quad и эта наименьшая величина, какъ видно изъ найденной нами формулы, стремится, при безпредѣльномъ возрастаніи $n$, къ предѣлу $\dfrac{ab}{(a+b)^2(a+b+1)}$, отличному отъ нуля.

Въ простѣйшемъ случаѣ, когда
\[
a = b = 1,
\]

% --- Page 149 ---

\noindent непримѣнимость закона большихъ чиселъ очевидна; такъ какъ въ этомъ случаѣ всѣ предположенія
\[
m = 0,\ 1,\ 2,\ 3\ \ldots\ n
\]
имѣютъ одну и туже вѣроятность $\dfrac{1}{n+1}$.

\begin{flushright}
16-го января 1907 года.
\end{flushright}

\medskip

\S\,5. Выводамъ \S\,2 можно придать значительно большую общность; а именно, вмѣсто числа появленій нѣкотораго событія можно разсматривать сумму величинъ, связанныхъ въ цѣпь такимъ образомъ, что, когда одна изъ нихъ получаетъ опредѣленное значеніе, слѣдующая за ней становится независимою отъ предшествующихъ ей. Пусть будетъ
\[
x_1,\, x_2,\, \ldots,\, x_k,\, x_{k+1},\, \ldots
\]
безконечный рядъ величинъ, связанныхъ такимъ образомъ, что $x_{k+1}$ при всякомъ $k$ не зависитъ отъ
\[
x_1,\, x_2,\, \ldots,\, x_{k-1},
\]
когда извѣстно значеніе $x_k$.

Пусть далѣе совокупность чиселъ
\[
\alpha,\, \beta,\, \gamma,\, \ldots
\]
представляетъ всѣ возможныя значенія любой изъ нашихъ величинъ, а система
\begin{align*}
&p_{\alpha,\,\alpha},\quad p_{\alpha,\,\beta},\quad p_{\alpha,\,\gamma},\quad \ldots \\
&p_{\beta,\,\alpha},\quad p_{\beta,\,\beta},\quad p_{\beta,\,\gamma},\quad \ldots \\
&\cdots\cdots\cdots\cdots\cdots
\end{align*}
представляетъ вѣроятности, при данномъ значеніи $x_k$, величинѣ $x_{k+1}$ имѣть опредѣленное значеніе, при чемъ первый

% --- Page 150 ---

значокъ у $p$ указываетъ данную величину $x_k$, а второй~---
предполагаемую величину $x_{k+1}$; напримѣръ при
\[
x_k = \beta
\]
вѣроятность равенства
\[
x_{k+1} = \gamma
\]
имѣетъ значеніе $p_{\beta,\gamma}$.

Эти вѣроятности мы предполагаемъ независящими отъ
значка $k$, чтобы не очень усложнять наши обозначенія и
разсужденія.

Пусть наконецъ
\[
p_{\alpha}^{(k)},\; p_{\beta}^{(k)},\; p_{\gamma}^{(k)},\;\ldots
\]
соотвѣтственно означаютъ вѣроятности для $x_k$ имѣть значенія
\[
\alpha,\; \beta,\; \gamma,\;\ldots
\]
пока всѣ величины
\[
x_1,\; x_2,\; x_3,\;\ldots\;,
\]
остаются неопредѣленными. Числа
\begin{align*}
& p_{\alpha,\,\alpha},\; p_{\alpha,\,\beta},\; p_{\alpha,\,\gamma},\;\ldots \\
& p_{\beta,\,\alpha},\; p_{\beta,\,\beta},\; p_{\beta,\,\gamma},\;\ldots
\end{align*}
\[
\cdots\cdots\cdots\cdots\cdots\cdots
\]
должны быть, конечно, положительными, кромѣ того они дол\-жны удовлетворять равенствамъ
\begin{align*}
& p_{\alpha,\,\alpha} + p_{\alpha,\,\beta} + p_{\alpha,\,\gamma} + \ldots\ldots = 1, \\
& p_{\beta,\,\alpha} + p_{\beta,\,\beta} + p_{\beta,\,\gamma} + \ldots\ldots = 1,
\end{align*}
\[
\cdots\cdots\cdots\cdots\cdots\cdots\cdots
\]
никакими другими условіями эти числа не ограничены.

% --- Page 151 ---

Что же касается чиселъ
\[
p_{\alpha}^{(k)},\; p_{\beta}^{(k)},\; p_{\gamma}^{(k)},\;\ldots\;,
\]
которыя, конечно, должны удовлетворять условію
\[
p_{\alpha}^{(k)} + p_{\beta}^{(k)} + p_{\gamma}^{(k)} + \ldots = 1,
\]
то ихъ нельзя задавать независимо для всѣхъ значеній $k$;
напротивъ по значеніямъ
\[
p_{\alpha}',\; p_{\beta}',\; p_{\gamma}',\;\ldots
\]
можно вычислять послѣдовательно всѣ величины
\[
p_{\alpha}^{(k)},\; p_{\beta}^{(k)},\; p_{\gamma}^{(k)},\;\ldots
\]
на основаніи простыхъ формулъ
\begin{align*}
& p_{\alpha}^{(k+1)} = p_{\alpha,\,\alpha}\, p_{\alpha}^{(k)} + p_{\beta,\,\alpha}\, p_{\beta}^{(k)} + \ldots \\
& p_{\beta}^{(k+1)} = p_{\alpha,\,\beta}\, p_{\alpha}^{(k)} + p_{\beta,\,\beta}\, p_{\beta}^{(k)} + \ldots
\end{align*}
\[
\cdots\cdots\cdots\cdots\cdots\cdots\cdots
\]

Обращаясь къ математическимъ ожиданіямъ нашихъ ве\-личинъ при различныхъ данныхъ, обозначимъ по прежнему
символомъ
\[
a_k
\]
математическое ожиданіе $x_k$, пока всѣ величины
\[
x_1,\; x_2,\;\ldots\;, x_k,\;\ldots
\]
остаются неопредѣленными, а символами
\[
A_{\alpha}^{(i)},\; A_{\beta}^{(i)},\; A_{\gamma}^{(i)},\;\ldots
\]
математическія ожиданія величины
\[
x_{k+i}
\]

% --- Page 152 ---

соотвѣтственно при
\[
x_k = \alpha,\; x_k = \beta,\; x_k = \gamma,\;\ldots\;.
\]

При такихъ обозначеніяхъ не трудно установить слѣду\-ющія равенства
\begin{align*}
a_k &= p_{\alpha}^{(k)} \alpha + p_{\beta}^{(k)} \beta + p_{\gamma}^{(k)} \gamma + \ldots \\[4pt]
a_{k+i} &= p_{\alpha}^{(k)} A_{\alpha}^{(i)} + p_{\beta}^{(k)} A_{\beta}^{(i)} + p_{\gamma}^{(k)} A_{\gamma}^{(i)} + \ldots \\[4pt]
A_{\alpha}^{(i)} &= p_{\alpha,\,\alpha}\, A_{\alpha}^{(i-1)} + p_{\alpha,\,\beta}\, A_{\beta}^{(i-1)} + p_{\alpha,\,\gamma}\, A_{\gamma}^{(i-1)} + \ldots \\[4pt]
A_{\beta}^{(i)} &= p_{\beta,\,\alpha}\, A_{\alpha}^{(i-1)} + p_{\beta,\,\beta}\, A_{\beta}^{(i-1)} + p_{\beta,\,\gamma}\, A_{\gamma}^{(i-1)} + \ldots
\end{align*}
\[
\cdots\cdots\cdots\cdots\cdots\cdots\cdots
\]

На основаніи этихъ равенствъ мы докажемъ, что всѣ величины
\[
a_{k+i},\; A_{\alpha}^{(i)},\; A_{\beta}^{(i)},\; A_{\gamma}^{(i)},\;\ldots\;,
\]
при безпредѣльномъ возрастаніи числа $i$, приближаются къ
одному предѣлу. Для этой цѣли прежде всего замѣтимъ, что
въ силу приведенныхъ нами равенствъ число $a_{k+i}$ лежитъ ме\-жду наибольшимъ и наименьшимъ изъ чиселъ
\[
A_{\alpha}^{(i)},\; A_{\beta}^{(i)},\; A_{\gamma}^{(i)},\;\ldots\;,
\]
а эти послѣднія числа всѣ заключаются между наибольшимъ
и наименьшимъ изъ чиселъ
\[
A_{\alpha}^{(i-1)},\; A_{\beta}^{(i-1)},\; A_{\gamma}^{(i-1)},\;\ldots\;.
\]
Разсматривая затѣмъ разность
\[
A_{\alpha}^{(i)} - A_{\beta}^{(i)},
\]
на основаніи тѣхъ же формулъ получаемъ

% --- Page 153 ---

\[
A_{\alpha}^{(i)} - A_{\beta}^{(i)} = (p_{\alpha,\,\alpha} - p_{\beta,\,\alpha})\, A_{\alpha}^{(i-1)} + (p_{\alpha,\,\beta} - p_{\beta,\,\beta})\, A_{\beta}^{(i-1)} + \ldots\ldots
\]
и такъ какъ сумма
\[
(p_{\alpha,\,\alpha} - p_{\beta,\,\alpha}) + (p_{\alpha,\,\beta} - p_{\beta,\,\beta}) + \ldots
\]
равна нулю, то, отдѣляя въ системѣ
\[
p_{\alpha,\,\alpha} - p_{\beta,\,\alpha},\; p_{\alpha,\,\beta} - p_{\beta,\,\beta},\;\ldots
\]
положительныя числа отъ отрицательныхъ и замѣняя послѣд\-нія ихъ числовыми величинами, мы получимъ двѣ совокупно\-сти положительныхъ чиселъ, образующихъ одинаковыя суммы.
Важно замѣтить также, что эти суммы меньше едини\-цы; ибо члены одной изъ нихъ меньше
\[
p_{\alpha,\,\alpha},\; p_{\alpha,\,\beta},\; p_{\alpha,\,\gamma},\;\ldots
\]
а члены другой меньше
\[
p_{\beta,\,\alpha},\; p_{\beta,\,\beta},\; p_{\beta,\,\gamma},\;\ldots\;.
\]

По этому, замѣняя одни изъ чиселъ
\[
A_{\alpha}^{(i-1)},\; A_{\beta}^{(i-1)},\; A_{\gamma}^{(i-1)},\;\ldots
\]
наибольшимъ, а другія наименьшимъ изъ нихъ и обозначая
символомъ
\[
\triangle^{(i-1)}
\]
разность между наибольшимъ и наименьшимъ изъ нашихъ
чиселъ
\[
A_{\alpha}^{(i-1)},\; A_{\beta}^{(i-1)},\; A_{\gamma}^{(i-1)},\;\ldots\;,
\]
мы можемъ установить неравенство

числ.\ знач.\ $(A_{\alpha}^{(i)} - A_{\beta}^{(i)}) < h\, \triangle^{(i-1)}$.

гдѣ $h$ равняется суммѣ всѣхъ положительныхъ чиселъ системы
\[
p_{\alpha,\,\alpha} - p_{\beta,\,\alpha},\; p_{\alpha,\,\beta} - p_{\beta,\,\beta},\;\ldots
\]

% --- Page 154 ---

и меньше единицы.

Совершенно подобное же заключеніе можно вывести от\-носительно разности любыхъ двухъ чиселъ системы
\[
A_{\alpha}^{(i)},\; A_{\beta}^{(i)},\; A_{\gamma}^{(i)},\;\ldots\ldots
\]
По этому, разсматривая вмѣсто
\[
A_{\alpha}^{(i)} - A_{\beta}^{(i)}
\]
разность $\triangle^{(i)}$ между наибольшимъ и наименьшимъ изъ чиселъ
\[
A_{\alpha}^{(i)},\; A_{\beta}^{(i)},\; A_{\gamma}^{(i)},\;\ldots\;,
\]
мы можемъ установить неравенство
\[
\triangle^{(i)} < H\,\triangle^{(i-1)}
\]
гдѣ $H$ означаетъ нѣкоторое постоянное число, лежащее меж\-ду~0 и~1.

Это неравенство показываетъ, что $\triangle^{(i)}$ приближается къ
предѣлу нуль, когда $i$ возрастаетъ безпредѣльно.

Слѣдовательно при безпредѣльномъ возрастаніи числа $i$
всѣ количества
\[
a_{k+i},\; A_{\alpha}^{(i)},\; A_{\beta}^{(i)},\; A_{\gamma}^{(i)},\;\ldots
\]
приближаются къ одному предѣлу, отъ котораго они отли\-чаются на величину меньшую $\triangle^{(i)}$; вмѣстѣ съ тѣмъ имѣемъ
\[
\triangle^{(i)} < C\, H^i,
\]
гдѣ $C$ и $H$ числа постоянныя и
\[
0 < H < 1.
\]

Обращаясь затѣмъ къ математическому ожиданію квадрата
\[
(x_1 - a_1 + x_2 - a_2 + \ldots + x_n - a_n)^2,
\]
мы разложимъ этотъ квадратъ на такія же слагаемыя какъ

% --- Page 155 ---

и въ \S\,2, полагая
\[
x_k - a_k = z_k.
\]
И въ силу доказаннаго, разсуждая совершенно также какъ въ
\S\,2, мы легко приходимъ къ неравенствамъ

мат.\ ожид.\ $z_k(z_1 + z_2 + \ldots + z_{k-1}) < D(H + H^2 + \ldots + H^{k-1})$

\noindent и

мат.\ ожид.\ $(x_1 - a_1 + x_2 - a_2 + \ldots + x_n - a_n)^2 < Gn,$

\noindent гдѣ $D$ и $G$ числа постоянныя.

Съ другой стороны, сравнивая математическое ожиданіе
\[
(x_1 - a_1 + x_2 - a_2 + \ldots + x_n - a_n)^2
\]
съ математическимъ ожиданіемъ
\[
(x_1 + x_2 + \ldots + x_n - na)^2,
\]
гдѣ
\[
a = \text{пред.}\; a_{k+i}\quad (i = \infty),
\]
находимъ, что разность математическихъ ожиданій этихъ
квадратовъ равна
\[
(a_1 - a + a_2 - a + \ldots + a_n - a)^2
\]
и сохраняетъ конечное значеніе, при безпредѣльномъ возра\-станіи числа $n$. Слѣдовательно при безпредѣльномъ возра\-станіи числа $n$ математическое ожиданіе квадрата
\[
\left(\frac{x_1 + x_2 + \ldots + x_n}{n} - a\right)^2
\]
должно приближаться къ предѣлу нуль.

А отсюда тотчасъ вытекаетъ для даннаго случая законъ
большихъ чиселъ: какъ бы малы ни были положительныя чи\-сла $\varepsilon$ и $\gamma$, вѣроятность выполненія неравенствъ
\[
-\varepsilon < \frac{x_1 + x_2 + \ldots + x_n}{n} - a < \varepsilon
\]

% --- Page 156 ---

будетъ больше $1 - \gamma$ для всѣхъ достаточно большихъ значе\-ній~$n$.

Итакъ независимость величинъ не составляетъ необхо\-димаго условія для существованія закона большихъ чиселъ.

\bigskip

\noindent\hspace{2em}25-го марта 1907 года

\begin{flushright}
\textbf{А. Марковъ.}
\end{flushright}

\end{document}
